% Work/MAT — revidovana hodina 30: pravidla pro odmocniny, bez racionalnich exponentu.
% Navaznost: 29 definice druhe a treti odmocniny; 31 racionalni exponent.
% Ucebnice: ucebnice/01_Zakladni_poznatky/Mocniny_a_odmocniny/Mocniny a Odmocniny.tex
% (Odmocniny, Caste chyby), Realisticky.cz 020712--020713 Pocitani s odmocninami.
% Zadani a reseni jsou nova; soucin/podil vzdy s presnymi podminkami.
\documentclass[11pt,a4paper]{article}
% Spolecny styl pro ucitelske pripravy MAT (Work/MAT)
% Verze 1.1, 2026-10-23
\NeedsTeXFormat{LaTeX2e}
\ProvidesPackage{priprava}[2026/10/23 v1.1 Sablona ucitelskych priprav MAT]

% Zakladni typografie a cestina
\RequirePackage[T1]{fontenc}
\RequirePackage[utf8]{inputenc}
\RequirePackage[czech]{babel}
\RequirePackage{lmodern}
\RequirePackage{microtype}
\RequirePackage{amsmath,amssymb}
\RequirePackage{xcolor}
\RequirePackage{enumitem}
\RequirePackage[most]{tcolorbox}
\RequirePackage{tabularx,booktabs,array}
\RequirePackage{fancyhdr}
\RequirePackage{needspace}
\RequirePackage{geometry}
\RequirePackage[hidelinks]{hyperref}

\geometry{a4paper,top=17mm,bottom=18mm,left=23mm,right=23mm}
\definecolor{notegray}{gray}{0.39}
\setlength{\parindent}{0pt}
\setlength{\parskip}{2.5pt plus 1pt minus 0.5pt}
\setlist[itemize]{topsep=2pt,itemsep=1pt,parsep=0pt,leftmargin=1.45em}
\setlist[enumerate]{topsep=2pt,itemsep=2pt,parsep=0pt,leftmargin=1.65em}
\setcounter{secnumdepth}{0}
\renewcommand{\arraystretch}{1.12}

% Zahlavi: vlevo rocnik | globalni tema, vpravo cislo: nazev hodiny
% Zapati zustava zcela prazdne.
\pagestyle{fancy}
\fancyhf{}
\renewcommand{\headrulewidth}{0.35pt}
\renewcommand{\footrulewidth}{0pt}
\setlength{\headheight}{14pt}
\newcommand{\nastavPripravu}[4]{%
  \fancyhead[L]{\small #1\enspace\textbar\enspace #2}%
  \fancyhead[R]{\small #3: #4}%
}

% Minimalisticke cernobile boxy. Kazda definice zvlast.
\tcbset{commonlesson/.style={enhanced,breakable=false,colback=white,colframe=black,
  colbacktitle=white,coltitle=black,boxrule=0.45pt,
  arc=0.8mm,outer arc=0.8mm,left=3mm,right=3mm,top=1.6mm,bottom=1.6mm,
  toptitle=1mm,bottomtitle=0.7mm,fonttitle=\bfseries,
  attach boxed title to top left={xshift=2.5mm,yshift=-1.15mm},
  boxed title style={colback=white,colframe=white,boxrule=0pt,left=1mm,right=1mm}}}
\newtcolorbox{zadaniBox}[1][]{commonlesson,title={Zadání},#1}
\newtcolorbox{prikladBox}[1][]{commonlesson,title={Příklad},#1}
\newtcolorbox{definiceBox}[1][]{commonlesson,title={Definice},#1}

% Spolecny box pro zadani a jeho reseni, bez mezery mezi castmi.
% Pouziti:
%   \begin{prikladReseniBox}
%     Zadani ulohy.
%     \tcblower
%     \begin{align*} ... \end{align*}
%   \end{prikladReseniBox}
% Spodni cast ma jemny sedy podklad a delici linku; lze v ni sazet
% vice radku matematiky. Obe casti maji jen jeden spolecny ramecek.
\tcbset{solutionlesson/.style={
  commonlesson,
  enhanced jigsaw,
  breakable,
  colbacklower=black!3,
  segmentation style={solid,draw=black!25,line width=0.35pt},
  fontlower=\small\color{notegray},
  before lower={\noindent\textbf{Řešení:}\enspace},
  leftlower=3mm,rightlower=3mm
}}
\newtcolorbox{prikladReseniBox}[1][]{solutionlesson,title={Příklad},#1}
\newtcolorbox{zadaniReseniBox}[1][]{solutionlesson,title={Zadání},#1}

% Pomucky pro autora pripravy: sedive poznamky; \reseni pro starsi pripravy
\newcommand{\poznamka}[1]{\par\begingroup\small\color{notegray}\itshape
  \textnormal{\textbf{Poznámka pro učitele.}}\enspace #1\par\endgroup}
\newcommand{\reseni}[1]{\par\begingroup\small\color{notegray}
  \textbf{Očekávané řešení:}\enspace #1\par\endgroup}
\newcommand{\blok}[1]{\needspace{5\baselineskip}\subsection*{#1}}
\newcommand{\vyrok}[1]{\uv{#1}}

\nastavPripravu{1. ročník}{Mocniny a odmocniny}{30}{Pravidla pro odmocniny}
\begin{document}
\poznamka{\textbf{45minutové jádro:} 4 min připomenutí druhé a třetí odmocniny;
9 min součin, podíl a Příklad 1; 8 min Samostatná práce 1 (a--d);
8 min vytýkání a vnášení činitele, Příklad 2; 7 min Samostatná práce 2 (a--d);
6 min sčítání podobných odmocnin a Příklad 3; 3 min shrnutí.
Další položky a všechny úlohy s proměnnou patří do rezervy.
\textbf{V této hodině nepoužívej zápis s racionálním exponentem.}}

\blok{Od významu odmocniny k pravidlům pro počítání}
V hodině 29 jsme určili $\sqrt{36}=6$ a $\sqrt[3]{-8}=-2$.
Dnes budeme odmocniny \emph{upravovat přesně}, bez desetinných aproximací.
Základem je skutečnost, že druhá odmocnina je \textbf{nezáporná},
kdežto třetí odmocnina smí být i záporná.

\blok{Společná tabulka pravidel pro odmocniny}
\begin{center}\small
\begin{tabularx}{\linewidth}{@{}>{\raggedright\arraybackslash}p{35mm}>{\raggedright\arraybackslash}p{58mm}X@{}}
\toprule
\textbf{Pravidlo} & \textbf{Zápis} & \textbf{Podmínky} \\ \midrule
Součin pod odmocninou & $\sqrt[n]{ab}=\sqrt[n]a\,\sqrt[n]b$ & $n\ge2$: pro sudé $n$ máme $a,b\ge0$; pro liché $n$ libovolná reálná $a,b$ \\
Podíl pod odmocninou & $\sqrt[n]{\dfrac ab}=\dfrac{\sqrt[n]a}{\sqrt[n]b}$ & pro sudé $n$: $a\ge0$, $b>0$; pro liché $n$: $a\in\mathbb R$, $b\ne0$ \\
Vyjmutí činitele & $\sqrt[n]{b^n a}=b\sqrt[n]a$ & pro sudé $n$: $b\ge0$, $a\ge0$; pro liché $n$: $a,b\in\mathbb R$ \\
Vnesení činitele & $b\sqrt[n]a=\sqrt[n]{b^n a}$ & stejné podmínky; u sudého $n$ a záporného $b$ viz poznámka níže \\
Podobné odmocniny & $m\sqrt a+k\sqrt a=(m+k)\sqrt a$ & $a\ge0$, $m,k\in\mathbb R$ \\
Mocnina odmocniny & $(\sqrt[n]a)^n=a$ & pro sudé $n$ jen $a\ge0$; pro liché $n$ libovolné $a$ \\
Odmocnina mocniny & $\sqrt[n]{b^n}=|b|$ pro sudé $n$, $=b$ pro liché $n$ & $b\in\mathbb R$, $n\ge2$ \\
\bottomrule
\end{tabularx}\normalsize\end{center}
\poznamka{Pro záporný činitel $b$ a sudé $n$ platí
$b\sqrt[n]a=-\sqrt[n]{b^n a}$, nikoliv pravidlo bez mínusu.
Příkladem je $-2\sqrt3=-\sqrt{12}$, zatímco $\sqrt{12}>0$.
\textbf{Nevytvářej z toho nové pravidlo:} rozhodující je nezápornost sudé odmocniny.}

\blok{Proč funguje součin a podíl?}
Pro $a,b\ge0$ je číslo $\sqrt a\,\sqrt b$ nezáporné a jeho druhá mocnina je
$(\sqrt a\,\sqrt b)^2=ab$. Musí tedy být rovno $\sqrt{ab}$.
Obdobně pro $b>0$ má $\sqrt a/\sqrt b$ druhou mocninu $a/b$.
U třetí odmocniny použijeme třetí mocninu; znaménko nás neomezuje.
\textbf{Neplatí však obecně} $\sqrt{a+b}=\sqrt a+\sqrt b$.

\begin{prikladReseniBox}[title={Příklad 1: součin a podíl odmocnin},breakable=false]
Uprav $\sqrt{9\cdot 16}$, $\sqrt{\dfrac{49}{81}}$, $\sqrt[3]{-8\cdot27}$ a
$\sqrt[3]{\dfrac{-125}{8}}$.
\tcblower
\[
\sqrt{9\cdot16}=3\cdot4=12,\quad
\sqrt{\frac{49}{81}}=\frac79,\quad
\sqrt[3]{-216}=-6,\quad \sqrt[3]{-125/8}=-\frac52.
\]
\end{prikladReseniBox}

\begin{zadaniReseniBox}[title={Samostatná práce 1: vhodný rozklad},breakable=false]
Uprav bez kalkulačky:
\begin{enumerate}[label=\alph*)]
\item $\sqrt{4\cdot 25}$; \item $\sqrt{\frac{36}{121}}$;
\item $\sqrt[3]{-27\cdot8}$; \item $\sqrt[3]{\frac{64}{-125}}$;
\item $\sqrt{12}\,\sqrt3$; \item $\frac{\sqrt{72}}{\sqrt2}$.
\end{enumerate}
\tcblower
\begin{enumerate}[label=\alph*)]
\item $10$; \item $6/11$; \item $-6$; \item $-4/5$;
\item $\sqrt{36}=6$; \item $\sqrt{36}=6$.
\end{enumerate}
\end{zadaniReseniBox}

\blok{Vynášení a vnášení činitele pod odmocnítko}
Pro nezáporné $b$ a $a$ lze zápis převádět oběma směry:
$\sqrt{b^2a}=b\sqrt a$ a $b\sqrt a=\sqrt{b^2a}$.
U třetí odmocniny může být $b$ i záporné.

\begin{prikladReseniBox}[title={Příklad 2: oba směry},breakable=false]
Uprav $\sqrt{75}$, $\sqrt[3]{-54}$, $4\sqrt3$ a $-2\sqrt3$.
\tcblower
$\sqrt{75}=\sqrt{25\cdot3}=5\sqrt3$; 
$\sqrt[3]{-54}=\sqrt[3]{-27\cdot2}=-3\sqrt[3]2$;
$4\sqrt3=\sqrt{16\cdot3}=\sqrt{48}$;
$-2\sqrt3=-\sqrt{4\cdot3}=-\sqrt{12}$.
\end{prikladReseniBox}

\begin{zadaniReseniBox}[title={Samostatná práce 2: vyjmi nebo vnes činitele},breakable=false]
\begin{enumerate}[label=\alph*)]
\item $\sqrt{98}$; \item $\sqrt[3]{-128}$;
\item $3\sqrt5$; \item $-4\sqrt2$;
\item $2\sqrt[3]{-7}$; \item $\sqrt{288}$.
\end{enumerate}
\tcblower
\begin{enumerate}[label=\alph*)]
\item $7\sqrt2$; \item $-4\sqrt[3]2$; \item $\sqrt{45}$;
\item $-\sqrt{32}$; \item $\sqrt[3]{-56}$; \item $12\sqrt2$.
\end{enumerate}
\end{zadaniReseniBox}

\blok{Sčítání podobných odmocnin}
Stejně jako $3x+5x=8x$ můžeme sečíst shodné odmocninové členy.
\begin{prikladReseniBox}[title={Příklad 3: nejprve uprav, potom sečti},breakable=false]
Zjednoduš $2\sqrt3+5\sqrt3-\sqrt3$ a $\sqrt{12}+\sqrt{27}$.
\tcblower
$2\sqrt3+5\sqrt3-\sqrt3=(2+5-1)\sqrt3=6\sqrt3$.
Dále $\sqrt{12}+\sqrt{27}=2\sqrt3+3\sqrt3=5\sqrt3$.
\end{prikladReseniBox}
\begin{zadaniReseniBox}[title={Samostatná práce 3: co lze a nelze sloučit?},breakable=false]
Zjednoduš:
\begin{enumerate}[label=\alph*)]
\item $3\sqrt2-7\sqrt2$; \item $2\sqrt{18}+\sqrt8$;
\item $4\sqrt5+\sqrt3$; \item $\sqrt{45}-\sqrt{20}$.
\end{enumerate}
\tcblower
\begin{enumerate}[label=\alph*)]
\item $-4\sqrt2$; \item $6\sqrt2+2\sqrt2=8\sqrt2$;
\item nelze sloučit; \item $3\sqrt5-2\sqrt5=\sqrt5$.
\end{enumerate}
\end{zadaniReseniBox}

\blok{Rozšiřující příklad 1: proměnná není automaticky kladná}
\begin{prikladReseniBox}[title={Rozšiřující příklad 1: absolutní hodnota},breakable=false]
Pro reálná $x,y$ uprav $\sqrt{36x^2}$, $\sqrt{(-5y)^2}$,
$\sqrt[3]{-8x^3}$ a vysvětli, kdy lze vynechat absolutní hodnotu.
\tcblower
$6|x|$, $5|y|$, $-2x$. U sudé odmocniny nezápornost vyžaduje
absolutní hodnotu; pokud předem víme $x\ge0$, pak $|x|=x$.
\end{prikladReseniBox}

\blok{Domácí úkol 1}
\begin{zadaniReseniBox}[title={Domácí úkol 1: smíšené odmocniny},breakable=false]
\begin{enumerate}[label=\alph*)]
\item $\sqrt{\frac{144}{25}}$; \item $\sqrt[3]{-250}$;
\item $\sqrt{200}$; \item $-3\sqrt7$ (vnes činitel);
\item $\sqrt{48}+\sqrt{75}$;
\item $5\sqrt2-\sqrt8+2\sqrt{18}$;
\item Rozhodni, zda $\sqrt{9+16}=\sqrt9+\sqrt{16}$; zdůvodni.
\end{enumerate}
\tcblower
\begin{enumerate}[label=\alph*)]
\item $12/5$; \item $-5\sqrt[3]2$; \item $10\sqrt2$;
\item $-\sqrt{63}$; \item $4\sqrt3+5\sqrt3=9\sqrt3$;
\item $5\sqrt2-2\sqrt2+6\sqrt2=9\sqrt2$;
\item ne: vlevo $5$, vpravo $7$.
\end{enumerate}
\end{zadaniReseniBox}
\poznamka{Na hodině 31 se teprve objeví zlomkový exponent.
Předcházející pravidla jsou o skutečné hodnotě odmocniny a musí
respektovat obor reálných čísel.}
\end{document}
